\documentclass[12pt]{article}
\usepackage[T1]{fontenc}
\usepackage[utf8]{inputenc}
\usepackage[brazil]{babel}
\usepackage[margin=1in]{geometry}
\usepackage{ragged2e}
\usepackage[normalem]{ulem}
\setlength{\parindent}{0pt}
\setlength{\parskip}{0.8\baselineskip}
\begin{document}
\justifying
\noindent Verifica-se que a questão objeto da controvérsia jurídica suscitada no recurso manejado pela parte recorrente esteve afetada no representativo da controvérsia - \textbf{Tema n. }\textbf{205} da TNU - (PU no processo n. PEDILEF 0500012-70.2015.4.05.8013/AL),  afetado em 25/02/2019 foi julgado (trânsito em julgado em 26/05/2020), por meio do qual foi elucidada a seguinte questão:

\noindent \textit{"Saber se é possível o enquadramento de atividade como especial por exposição a agentes biológicos, quando os serviços prestados não são aqueles descritos no Anexo IV do Decreto 3.048/99."}

\noindent Restou firmada a seguinte tese:

\noindent \textit{"a) para reconhecimento da natureza especial de tempo laborado em exposição a agentes biológicos não é necessário o desenvolvimento de uma das atividades arroladas nos Decretos de regência, sendo referido rol meramente exemplificativo; b) entretanto, é necessária a comprovação em concreto do risco de exposição a microorganismos ou parasitas infectocontagiosos, ou ainda suas toxinas, em medida denotativa de que o risco de contaminação em seu ambiente de trabalho era superior ao risco em geral, devendo, ainda, ser avaliado, de acordo com a profissiografia, se tal exposição tem um caráter indissociável da produção do bem ou da prestação do serviço, independentemente de tempo mínimo de exposição durante a jornada (Tema 211/TNU)." (Tema 205 TNU).}

\noindent Com efeito, depreende-se da leitura do voto condutor do acórdão que o mesmo se encontra em conformidade com o decidido no  tema pela TNU.

\noindent A proposito, consta do acordao hostilizado: (...) Convalidação de Estudos e Reconhecimento de Diploma, ACESSO, DIREITO À EDUCAÇÃO

\noindent Nessas condições, o conteúdo do Acórdão recorrido é consentâneo com o entendimento da TNU, firmado em julgamento de recurso representativo de controvérsia, o que atrai a incidência da regra prevista pelo art. 14, III, \textit{b,} da Resolução CJF n. 586/2019:

\noindent \textit{\textbf{Art. 14.}}\textit{ Decorrido o prazo para contrarrazões, os autos serão conclusos ao magistrado responsável pelo exame preliminar de admissibilidade, que deverá, de forma sucessiva: (...) }\textit{\textbf{III }}\textit{- negar seguimento a pedido de uniformização de interpretação de lei federal interposto contra acórdão que esteja em conformidade com entendimento consolidado: }\textit{\textbf{b)}}\textit{ em }\uline{\textit{\textbf{recurso representativo de controvérsia pela Turma Nacional de Uniformização}}}\textit{ ou em pedido de uniformização de interpretação de lei dirigido ao Superior Tribunal de Justiça;(...).}

\noindent Posto isso, com arrimo no \uline{\textbf{art. 14, III, }}\uline{\textit{\textbf{b}}}\textit{,} da Resolução CJF n. 586/2019 c/c art. 1.030, I, do CPC, \textbf{NEGO SEGUIMENTO }\textbf{a}o\textbf{ PU}\textbf{.}

\noindent Intimem-se. Aguarde-se o decurso do prazo recursal. Após, certifique-se o trânsito em julgado e devolva-se ao juízo de origem.
\end{document}
